Todo.

 1.) just like the test function for the velocity comes from a space that is a closure of a continuous space, the concentration test function should come from a continuous one as well (!!!!!!!!) 
 
 2.) generalized model: update every part (wf, ieq, theorem) accordingly with the new theta part and the new part_2, part_3 C part.

 3.)  add boundary conditions to the H and A set of equations and add that to the list in the main theorem

 4.) $\mathbf{u} C$ product to be in $L^2$ in the hydrostatic weak formulation - spaces.  

 5.) Officially the Trace theorem can be used with $\norm{u}_{W^{1,p}(U)}$ if $u$ is in that space, but in the inland case, $C^\epsilon$ is not a priori in the $H^1$ in space, we need to say something here when we pass from the line integral to the main domain.
  
Future Work.

  \item Use more sophisticated ground-level boundary conditions instead of $\partial_3 C = 0$ in the inland case, for example \cite{monin1959boundary}.
  
  \item We could consider multiple contaminants.
  
  \item It is possible and not unreasonable in certain cases to use isotropic diffusivities $K_x = K_y = K_z \neq 0$.
  
  \item It is an option to simply take $K_x = 0$, without using $\epsilon$ as suggested by \cite{hosseini2013}, the advection/convection effect being dominant compared to diffusion. It could also be possible to assume that $K_x = K_z$ due to symmetry in some cases(p16).
  
  \item Try to apply the
  \[ \frac{\partial \Phi}{\partial t} + u \frac{\partial \Phi}{\partial n} = 0 \]
  radiation condition derived by the wave equation on the lateral boundary, where $u$ is the current in the direction to the normal to the boundary. This is suggested articles such as \cite{marsaleix2006considerations} and \cite{wang2014ocean}.